\documentclass{article}

% packages
\usepackage{amsmath, amsthm, thmtools, amsfonts, amssymb, luacode, catchfile, tikzducks, hyperref, ifthen}
\ifcsname c@kobocompile\endcsname
	\usepackage[a5paper, total={1072pt, 1448pt}, margin=10pt, includeheadfoot]{geometry} % set page margins
\else
	\usepackage[a4paper, margin=50pt, includeheadfoot]{geometry}
\fi
\usepackage[shortlabels]{enumitem}
\usepackage[skip=3pt, indent=0pt]{parskip}

% language
\usepackage[bidi=basic, layout=tabular, provide=*]{babel}
\ifcsname c@english\endcsname
	\babelprovide[main, import]{english}
\else
	\babelprovide[main, import]{hebrew}
	\babelprovide{rl}
\fi
%\babelfont{rm}{Libertinus Serif}
\babelfont{rm}[Renderer=Harfbuzz]{Libertinus Serif}
\babelfont{sf}{Libertinus Sans}
\babelfont{tt}{Libertinus Mono}

% style
\AddToHook{cmd/section/before}{\clearpage}	% Add line break before section
\linespread{1.3}
\setcounter{secnumdepth}{0}		% Remove default number tags from sections, this won't do well with theorems
\AtBeginDocument{\setlength{\belowdisplayskip}{3pt}}
\AtBeginDocument{\setlength{\abovedisplayskip}{3pt}}
\graphicspath{ {../images/} }

% operators
\DeclareMathOperator\cis{cis}
\DeclareMathOperator\Sp{Sp}
\DeclareMathOperator\tr{tr}
\DeclareMathOperator\im{Im}
\DeclareMathOperator\re{Re}
\DeclareMathOperator\diag{diag}
\DeclareMathOperator*\lowlim{\underline{lim}}
\DeclareMathOperator*\uplim{\overline{lim}}
\DeclareMathOperator\rng{rng}
\DeclareMathOperator\Sym{Sym}
\DeclareMathOperator\Arg{Arg}
\DeclareMathOperator\Log{Log}
\DeclareMathOperator\dom{dom}
\DeclareMathOperator\supp{Supp}
\DeclareMathOperator\var{Var}
\DeclareMathOperator\cov{Cov}

% commands
%\renewcommand\qedsymbol{\textbf{מש''ל}}
%\renewcommand\qedsymbol{\fbox{\emoji{lizard}}}
\newcommand{\Aa}[0]{\mathcal{A}}
\newcommand{\Bb}[0]{\mathcal{B}}
\newcommand{\CC}[0]{\mathbb{C}}
\newcommand{\Cc}[0]{\mathcal{C}}
\newcommand{\EE}[0]{\mathbb{E}}
\newcommand{\FF}[0]{\mathbb{F}}
\newcommand{\Ff}[0]{\mathcal{F}}
\newcommand{\Ii}[0]{\mathcal{I}}
\newcommand{\Gg}[0]{\mathcal{G}}
\newcommand{\Ll}[0]{\mathcal{L}}
\newcommand{\Mm}[0]{\mathcal{M}}
\newcommand{\NN}[0]{\mathbb{N}}
\newcommand{\Nn}[0]{\mathcal{N}}
\newcommand{\PP}[0]{\mathbb{P}}
\newcommand{\Pp}[0]{\mathcal{P}}
\newcommand{\QQ}[0]{\mathbb{Q}}
\newcommand{\RR}[0]{\mathbb{R}}
\newcommand{\Rr}[0]{\mathcal{R}}
\newcommand{\Ss}[0]{\mathcal{S}}
\newcommand{\TT}[0]{\mathbb{T}}
\newcommand{\Uu}[0]{\mathcal{U}}
\newcommand{\Vv}[0]{\mathcal{V}}
\newcommand{\Ww}[0]{\mathcal{W}}
\newcommand{\ZZ}[0]{\mathbb{Z}}
\newcommand{\acts}[0]{\circlearrowright}
\newcommand{\explain}[2] {
	\begin{flalign*}
		 && \text{#2} && \text{#1}
	\end{flalign*}
}
\newcommand{\maketitleprint}[0]{ \begin{center}
	%\begin{tikzpicture}[scale=3]
	%	\duck[graduate=gray!20!black, tassel=red!70!black]
	%\end{tikzpicture}	
	\includegraphics[width=6cm]{cover}
\end{center}
}

% theorem commands
\newtheoremstyle{c_remark}
	{}	% Space above
	{}	% Space below
	{}% Body font
	{}	% Indent amount
	{\bfseries}	% Theorem head font
	{}	% Punctuation after theorem head
	{.5em}	% Space after theorem head
	{\thmname{#1}\thmnumber{ #2}\thmnote{ \normalfont{\text{(#3)}}}}	% head content
\newtheoremstyle{c_definition}
	{3pt}	% Space above
	{3pt}	% Space below
	{}% Body font
	{}	% Indent amount
	{\bfseries}	% Theorem head font
	{}	% Punctuation after theorem head
	{.5em}	% Space after theorem head
	{\thmname{#1}\thmnumber{ #2}\thmnote{ \normalfont{\text{(#3)}}}}	% head content
\newtheoremstyle{c_plain}
	{3pt}	% Space above
	{3pt}	% Space below
	{\itshape}% Body font
	{}	% Indent amount
	{\bfseries}	% Theorem head font
	{}	% Punctuation after theorem head
	{.5em}	% Space after theorem head
	{\thmname{#1}\thmnumber{ #2}\thmnote{ \text{(#3)}}}	% head content

\ifcsname c@english\endcsname
	\theoremstyle{plain}
	\newtheorem{theorem}{Theorem}[section]
	\newtheorem{lemma}[theorem]{Lemma}
	\newtheorem{proposition}[theorem]{Proposition}
	\newtheorem*{proposition*}{Proposition}
	%\newtheorem{corollary}[theorem]{אין חלופה עברית}

	\theoremstyle{definition}
	\newtheorem{definition}[theorem]{Definition}
	\newtheorem*{definition*}{Definition}
	\newtheorem{example}{Example}[section]
	\newtheorem{exercise}{Exercise}[section]

	\theoremstyle{remark}
	\newtheorem*{remark}{Remark}
	\newtheorem*{solution}{Solution}
	\newtheorem{conclusion}[theorem]{Conclusion}
	\newtheorem{notation}[theorem]{Notation}
\else
	\theoremstyle{c_plain}
	\newtheorem{theorem}{משפט}[section]
	\newtheorem{lemma}[theorem]{למה}
	\newtheorem{proposition}[theorem]{טענה}
	\newtheorem*{proposition*}{טענה}
	%\newtheorem{corollary}[theorem]{אין חלופה עברית}

	\theoremstyle{c_definition}
	\newtheorem{definition}[theorem]{הגדרה}
	\newtheorem*{definition*}{הגדרה}
	\newtheorem{example}{דוגמה}[section]
	\newtheorem{exercise}{תרגיל}[section]

	\theoremstyle{c_remark}
	\newtheorem*{remark}{הערה}
	\newtheorem*{solution}{פתרון}
	\newtheorem{conclusion}[theorem]{מסקנה}
	\newtheorem{notation}[theorem]{סימון}
\fi

% Questions related commands
\newcounter{question}
\setcounter{question}{1}
\newcounter{sub_question}
\setcounter{sub_question}{1}

\ifcsname c@english\endcsname
	\newcommand{\question}[1][0]{
		\ifthenelse{#1 = 0}{}{\setcounter{question}{#1}}
		\section{Question \arabic{question}}
		\addtocounter{question}{1}
		\setcounter{sub_question}{1}
	}

	\newcommand{\subquestion}[1][0]{
		\ifthenelse{#1 = 0}{}{\setcounter{sub_question}{#1}}
		\subsection{Part \alph{sub_question}}
		\addtocounter{sub_question}{1}
	}
\else
	\newcommand{\question}[1][0]{
		\ifthenelse{#1 = 0}{}{\setcounter{question}{#1}}
		\section{שאלה \arabic{question}}
		\addtocounter{question}{1}
		\setcounter{sub_question}{1}
	}

	\newcommand{\subquestion}[1][0]{
		\ifthenelse{#1 = 0}{}{\setcounter{sub_question}{#1}}
		\subsection{סעיף \localecounter{letters.gershayim}{sub_question}}
		\addtocounter{sub_question}{1}
	}
\fi

% import lua and start of document
\directlua{common = require ('../common')}

\GetEnv{AUTHOR}

% headers
\author{\AUTHOR}
\date\today

\title{פתרון מטלה 04 --- מבוא לטופולוגיה, 80516}
% chktex-file 9
% chktex-file 17

\begin{document}
\maketitle
\maketitleprint{}

\question{}
יהי $(X, \rho)$ מרחב מטרי.
נוכיח כי התנאים הבאים שקולים,
\begin{enumerate}
	\item $X$ מקיים את אקסיומת המנייה השנייה
	\item $X$ הוא מרחב לינדלוף
	\item $X$ ספרבילי
\end{enumerate}
\begin{proof}
	$1 \implies 2$.
	אקסיומת המנייה גוררת את כל אקסיומות המנייה האחרות, לכן בפרט $X$ מרחב לינדלוף.

	$2 \implies 3$.
	נרצה למצוא קבוצה צפופה בת־מנייה ב־$X$.
	נוכל לכסות את המרחב על־ידי כדורי $\epsilon > 0$ לכל נקודה במרחב, לכן בפרט יש כיסוי בן־מניה של כל המרחב על־ידי קבוצות פתוחות לכל $\epsilon$ כזה.
	נניח ש־${\{ B(x_n^m, \epsilon_m) \}}_{n = 1}^\infty$ קבוצה כזו עבור $\epsilon_m = \frac{1}{m}$ לכל $m$.
	נגדיר את הקבוצה $A = \{ x_n^m \mid n, m \in \NN \}$, זוהי כמובן קבוצה בת־מניה, ונטען ש־$A$ צפופה ב־$X$.
	נניח ש־$x \in X$, אם $x \in A$ אז סיימנו, אחרת לכל $m \in \NN$ קיים $n(m) \in \NN$ כך ש־$x \in B(x_{n(m)}^m, \frac{1}{m})$ מהגדרת הכיסויים, ולכן הסדרה ${( x_{n(m)}^m )}_{m = 1}^\infty$ מתכנסת ל־$x$ ולכן $x \in \overline{A}$.
	נסיק אם כך שאכן $A$ צפופה ובת־מניה ומעידה כי $X$ ספרבילי.

	$3 \implies 1$.
	נגדיר את הקבוצה $\Bb = \{ B(a, \frac{1}{n}) \mid a \in A, n \in \NN \}$ עבור $A$ קבוצה צפופה ב־$X$ שקיימת מספרביליות.
	אנו טוענים כי $\Bb$ בסיס בן־מניה ל־$X$.
	$|\Bb| = |\NN \times A| = \aleph_0$, ולכן נשאר להראות כי זהו אכן בסיס לטופולוגיה של $X$, אך למעשה זו טענה שהוכחה בהרצאה, ולכן נוכל להסיק ש־$X$ מקיימת את אקסיומת המנייה השנייה.
\end{proof}

\question{}
\subquestion{}
יהי $X$ מרחב מנייה שנייה ויהי $A \subseteq X$,
נראה שגם $A$ מרחב מנייה שנייה.
\begin{proof}
	יהי $\Bb$ בסיס בן־מניה של $X$,
	ויהי $\Bb_0 = \{ B \cap A \mid B \in \Bb \}$, זוהי קבוצה בת־מנייה, ואנו טוענים שאף בסיס של $A$.
	נניח ש־$U \subseteq A$ קבוצה פתוחה, אז קיימת $V \subseteq X$ כך ש־$U = V \cap A$ וכן $V$ פתוחה ב־$X$.
	בהתאם $V = \bigcup_{\alpha \in I} V_{\alpha}$ עבור $I$ קבוצת אינדקסים כך ש־$V_{\alpha} \in \Bb$ לכל $\alpha \in I$.
	לכן נובע ש־$U = (\bigcup_{\alpha \in I} V_{\alpha}) \cap A = \bigcup_{\alpha \in I} V_{\alpha} \cap A = \bigcup_{\alpha \in I} U_{\alpha}$ עבור $U_{\alpha} \in \Bb_0$.
	נסיק אם כך שאכן $\Bb_0$ בסיס לטופולוגיה של $A$, וראינו כבר כי הוא בן מנייה, ולכן $A$ מרחב מנייה שנייה.
\end{proof}

\subquestion{}
נבחן את המרחב המטרי $X = {[0, 1]}^\NN$ עם המטריקה $\rho(x, y) = \sup \{ |x(n) - y(n)| \mid n \in \NN \}$. \\
נראה שהוא אינו ספרבילי, אינו לינדלוף ואינו מנייה שנייה.
\begin{proof}
	נבחן את הקבוצה $X \supseteq A = \{ f \in X \mid \forall n, f(n) \in \{0, 1\} \}$, כלומר הסדרות שתמונתן חלקית ל־$\{0, 1\}$.
	לכל $x, y \in A$, מתקיים $\rho(x, y) = 0 \iff x = y$, מהגדרת המטריקה אך גם $\rho(x, y) = 1 \iff x \ne y$, כלומר זוהי המטריקה הדיסקרטית, ולכן היא משרה טופולוגיה דיסקרטית.

	נניח בשלילה ש־$X$ מרחב מנייה שנייה (ולכן שקול לשאר ההגדרות כנביעה משאלה 1), ולכן מסעיף א' נובע שגם $A$ מרחב מנייה שנייה.
	לכן קיים בסיס בן־מנייה ל־$A$, אבל כל יחידון צריך להופיע בבסיס, זאת שכן כל איבר בטופולוגיה הוא איחוד של איברי הבסיס, לכן נסיק ש־$A$ הוא הבסיס היחיד של $A$.
	נבחין כי $|A| = 2^\NN > \aleph_0$, ולכן ל־$A$ אין בסיס בן־מנייה בסתירה להיותו מרחב מנייה שנייה.
	נסיק שאכן $X$ איננו מרחב מנייה שנייה, אינו לינדלוף, ואינו ספרבילי.
\end{proof}

\question{}
נגדיר טופולוגיה על $\NN^2$ כך ש־$U$ פתוחה אם ורק אם $(0, 0) \notin U$ או שקיים $n$ כך שלכל $m > n$ הקבוצה $\{ k \in \NN \mid (m, k) \notin U \}$ היא סופית. \\
מרחב זה נקרא מרחב ארנס־פורט, נראה כי הוא לא מרחב מנייה שנייה.
\begin{proof}
	תהי $M \subseteq \NN$, נגדיר את הקבוצה $A = \{ (m, k) \in \NN^2 \mid m \in M, k \in \NN \} \cup \{ (m, k + 1) \in \NN^2 \mid m \notin M, k \in \NN \}$.
	לכל $m \in \NN$ הקבוצה $\{ k \in \NN \mid (m, k) \notin A \} \in \{ \emptyset, \{ 0 \} \}$ ובפרט סופית ולכן $A$ קבוצה פתוחה.
	נגדיר את הפונקציה $f : \Pp(\NN) \to \tau$ עבור $\tau$ הטופולוגיה שהגדרנו על־ידי מיפוי כל קבוצה $M$ לקבוצה $A$ המתאימה לה.
	זוהי פונקציה חד־חד ערכית, זאת שכן אם $M_0 \ne M_1$ אז ללא הגבלת הכלליות $m \in M_0 \setminus M_1$ ואז $(m, 0) \in f(M_0)$ אבל מההגדרה $(m, 0) \notin f(M_1)$.
	נסיק ש־$\aleph_0 < |\Pp(\NN)| \le |\tau|$.

	נניח ש־$\Bb$ בסיס בן־מנייה לנקודה $(0, 0)$.
	לכל $f(M)$ יש קבוצה פתוחה $U \in \Bb$ כך ש־$(0, 0) \in U \subseteq f(M)$.
	$\Bb$ בת־מנייה ולכן קיימת $U$ כך שיש אינסוף ערכי $f(M)$ המכילים את $U$.
	כלומר אם נקבע $A = f(M)$ ונגדיר את הסדרה $A_0 = A, A_1 = A_0 \setminus \{ (1, k) \mid k \in \NN \}, A_2 = \ldots$ אז נקבל ש־$U \subseteq A_i$ לכל $i \in \NN$, בפרט $U \subseteq \bigcap_{i \in \NN} A_i$.
	אבל $\bigcap_{i \in \NN} A_i$ לא מקיימת את התנאי של הסופיות ובפרט לא קבוצה פתוחה, נסיק ש־$U$ אף היא לא פתוחה, בסתירה ל־$U \in \Bb$.
	נסיק שמרחב זה הוא לא מנייה שנייה.
\end{proof}

\question{}
נוכיח שמכפלה סופית של מרחבים ספרביליים היא ספרבילית.
\begin{proof}
	נניח ש־$X, Y$ מרחבים ספרביליים ונבחן את $Z = X \times Y$.
	נניח גם ש־$A \subseteq X, B \subseteq Y$ קבוצות צפופות בנות־מניה המעידות על הספרביליות.
	נגדיר $C = A \times B$, זוהי קבוצה בת־מניה כמכפלה סופית של קבוצות בנות־מניה, אנו נראה שהיא צפופה ב־$Z$.
	ניעזר בטענה שקבוצה היא צפופה אם ורק אם היא חותכת כל קבוצה פתוחה במרחב, ונניח ש־$U = U_X \times U_Y \subseteq Z$ פתוחה, לכן מהגדרת מכפלות סופיות למרחבים טופולוגיים $U_X, U_Y$ פתוחות ב־$X, Y$.
	בהתאם $U_X \cap A, U_Y \cap B$ קבוצות לא ריקות, נניח ש־$x \in X, y \in Y$ איברים המעידים על כך.
	אז $C \cap U = (A \times B) \cap (U_X \cap U_Y) = (A \cap U_X) \times (B \cap U_Y) \ni (x, y)$, כלומר החיתוך של כל קבוצה פתוחה עם $C$ לא ריק, ולכן היא צפופה ב־$Z$.
\end{proof}

\question[6]
נוכיח שטופולוגיית הקופסה על $\RR^\NN$ אינה קשירה.
\begin{proof}
	נגדיר $A = \{ f : \NN \to \RR \mid f(n) \le 0 \}$ וכן $B = \{ f : \NN \to \RR \mid f(n) > 0 \}$.
	לכל $n \in \NN$, נבחין כי $A_n = \{ f(n) \mid f \in A \}$ היא הקבוצה $(-\infty, 0]$, 

	TODO: Complete.
\end{proof}

\question{}
\subquestion{}
נמצא דוגמה למרחב קשיר מקומית שאיננו קשיר.
\begin{solution}
	נגדיר $X = \{0, 1\}$ עם הטופולוגיה הדיסקרטית.
	המרחב איננו קשיר שכן $X = \{ 0 \} \cup \{ 1 \}$ ובמרחב דיסקרטי כל יחידון הוא פתוח.
	מן הצד השני המרחב קשיר מקומית ב־$0$, זאת שכן כל קבוצה כך ש־$A \subseteq X$ וגם $0 \in A$, אז $\{ 0 \} \subseteq A$ וכן $\{ 0 \}$ פתוחה.
\end{solution}

\subquestion{}
לכל $0 < n \in \NN$ נסמן $l_n = \{ (t, \frac{t}{n}) \mid t \in [0, 1] \}$ ונסמן $l = [0, 1] \times \{ 0 \}$.
נגדיר את המרחב $X = l \cup \bigcup_{n > 0} l_n$ ונראה שהוא קשיר אך לא קשיר מקומית.
\begin{proof}
	נבחין כי לכל $n \in \NN$ הקבוצה $l_n$ היא קטע ולכן קשירה כמסקנה מהתרגול.
	גם $l$ הוא קטע ולכן גם הוא קשיר.
	נבחין כי לכל $n \in \NN$, מתקיים $(0, 0) \in l_n$ וכן ש־$(0, 0) \in l$, ולכן מלמת כוכב המרחב הוא קשיר (ואף קשיר מסילתית).

	נראה שהמרחב איננו קשיר מקומית.
	נבחן את הקבוצה הפתוחה $U = B((\frac{1}{2}, \frac{1}{2}), \frac{1}{4}) \cap X$, זוהי קבוצה פתוחה כצמצום קבוצה פתוחה במרחב המטרי $\RR^2$.
	מהצד השני נבחין כי $U = \{ (t, t) \mid t \in (\frac{1}{3}, \frac{2}{3}) \} \cup \{ (t, \frac{t}{2}) \mid t \in (\frac{1}{2}, \frac{7}{10}) \}$.
	כלומר $U$ הוא איחוד קטעים זרים ופתוחים ולכן לא קשיר מקומית.
\end{proof}

\end{document}
